\section{Investigation and Source-Content Ranking}
Outcome review starts from the user's permitted directories and actual final changes. C1 accepts the allowed directory list, checks recorded and inferred file writers, and separately compares baseline hashes with final workspace hashes. It builds without final-check reconciliation: an unexplained final change cannot invent a writer. Failed and invalid calls remain diagnostics. Shell-derived writes remain inferences, while delegated task boundaries and actual child edits are counted separately.

The query command follows ancestors or descendants of a selected call and retains original arguments, returns and evidence edges. A child query also exposes the parent task boundary. The ancestor fraction is the number of ancestor nodes excluding the target divided by all nodes in the same session graph excluding the target. It measures graph coverage, not files or analyst time. Parent and child fractions are separate.

C2 ranks only nonempty tool returns in earlier assistant turns of the same native session; a return ending after the target starts is excluded. Same-turn calls, future returns and user messages are not candidates. Child ranking does not assume access to parent history. Direct read/skill returns with one file identity are grouped by path; shell and otherwise unmapped returns remain individual candidates. Repeated reads of a file take the maximum return score rather than accumulating points.

For an edit, raw old/new replacement text determines added token occurrences and added or replaced full lines. Tokens are signed numbers, identifiers or paths. Positive differences in token counts produce distinct added-token features; unchanged context adds no token points. Full new-line features retain any unchanged context on that changed line. For a write, exact previous content is required; unavailable content is reported instead of treating an overwrite as wholly new.

Rule v1 awards one point for the complete target path, one per distinct added token, and one per full added line found in the return. File scores take the maximum over returns. A tie interval is one plus the count of higher-scoring candidates through the count of candidates with at least that score. Stable display order does not break ties. The algorithm receives no scenario label or known injection path; the acceptance script looks up that path only after ranking.

These scores retrieve content-similar candidates. They do not establish model-internal influence. In particular, the wrong amounts also appear in real failing-test output, and 5.50 already occurs in a legitimate test input. We make no source-attribution, ranking-improvement or time-saving claim. The same-input comparison below retains this ambiguity.


\subsection{C3: exploratory operation rarity}
A frozen rule uses normal runs 030, 033, 034, 037 and 038 as references. A signature consists of tool, relation and normalized path, without node, version or session identifiers. CALL denotes a native invocation with null path, including failed or invalid calls. READ, WRITE and DELETE use recorded or inferred file edges built without final reconciliation. Parent delegated file summaries are excluded; actual child operations belong to their complete run. Failure does not manufacture file relations.

Let $k_s$ count reference runs containing signature $s$, once per complete run regardless of repetitions, and let $N$ be the number of reference runs. We use
\[
 p_s=\frac{k_s+1}{N+2},\qquad R_s=-\log_2 p_s.
\]
This is independent binary occurrence smoothing per signature, not a categorical distribution over paths; probabilities need not sum to one. Any unseen path or tool has $k_s=0$ and finite rarity. With five references, an unseen signature scores 2.807 bits and a signature in all five scores 0.222 bits. These are descriptive rarity scores, not attack probabilities.

Each native call takes the maximum over its signatures; CALL and maximum file scores remain separately inspectable. An empty file neighborhood has null file score. Scores are not added across edges, repeated operations or child sessions. Calls are ranked within one complete run with explicit tie intervals. No alert threshold or C1/C2 score combination is fitted.

Normal scores under the full reference are in-sample descriptions. Five additional folds each remove one normal run and score that held-out run and all five injection runs with the same four-reference model. The five folds contain five distinct held-out normal cases, not 25 new injection experiments. Four-reference and five-reference raw scores have different calibration. This is offline analysis, not a detector deployed in collection-time order.


\subsection{C4: a graph-independent flat-log baseline}
Both methods receive the same ten runs, native parent/child sessions, initial hashes, final files and authorized directories. The automated log baseline scans calls and maintains a path set; it does not import graph methods or read graph artifacts. It includes actual child calls and retains failures. The two methods share parsing, normalization, shell inference, authorization and the fixed C2-v1 content rule. File operations and source identities are obtained separately. This compares evidence organization under shared observation, not added system-call visibility.

Detection is checked by run and native write operation (session, call and path), excluding parent summaries. Source coverage includes only confirmed violating starts with full prior payload exposure in their own session. Manual labels and the known controlled path are consulted only after computation. We report candidate counts, tied ranks and the ancestor fraction with the target excluded from numerator and denominator. For a common volume unit we also compare graph-ancestor call identities with the flat log's same-session temporal prefix. Parent context is separate from child denominators. No analyst timing or tuned fusion with C3 is performed.
